\section{Importancia del proyecto en base al contexto y la evidencia}

\subsection{Breve contexto de la Reinserción en Chile}
La literatura especializada y la evidencia administrativa coinciden en que el empleo formal constituye el principal factor protector contra la reincidencia delictiva. En Chile, la reinserción social se concibe como un proceso sistemático e interinstitucional en el que convergen el Estado, la sociedad civil y el sector privado con el objetivo de integrar a personas condenadas e infractoras de ley, reducir los índices de reincidencia y fomentar conductas prosociales.

La estructura institucional de este sistema reposa sobre tres pilares fundamentales:
\begin{itemize}
    \item \textbf{Servicio Nacional de Reinserción Social Juvenil (SNRSJ):} Creado mediante la Ley N° 21.527 para reemplazar de forma gradual al Servicio Nacional de Menores (SENAME) en la atención de jóvenes de entre 14 y 18 años, operando bajo un modelo descentralizado, altamente especializado y con enfoque garantista \parencite{Ley21527CreaElServicioNacionalDeReinsercionJuvenil}.
    \item \textbf{Gendarmería de Chile:} Organismo responsable de ejecutar programas de intervención psicosocial criminológica y de empleabilidad para adultos, mediante capacitación en oficios, alianzas de intermediación laboral (como el programa +R) y la comercialización de productos fabricados en recintos penitenciarios a través de plataformas como REHACE \parencite{EstadisticaGeneralPenitenciaria2026}.
    \item \textbf{Articulación Público-Privada:} Redes de colaboración estratégica como \textit{Juntos por la Reinserción} (JxR), que agrupan a fundaciones, gremios y empresas para facilitar procesos sostenibles de capacitación y colocación laboral para personas con antecedentes penales.
\end{itemize}

Sin embargo, las personas que egresan del sistema penitenciario enfrentan un entorno de contratación severamente restringido por estigmas culturales, exigencias legales (como la presentación del certificado de antecedentes penales) y barreras procedimentales en las áreas de gestión de personas. En el caso de las mujeres privadas de libertad, la brecha de inserción se intensifica por dinámicas de doble penalización: la estigmatización social asociada al quebrantamiento de la norma de género y las severas cargas no remuneradas de cuidado familiar \parencite{CaracterizacionPersonasPrivadas2026}.

\subsection{Población Atendida y Caracterización del Sistema Penitenciario}
Según la Estadística General Penitenciaria de Gendarmería de Chile, el sistema atiende a un total de 162.043 personas a nivel nacional, compuestas por un 88,2\% de hombres y un 11,8\% de mujeres \parencite{EstadisticaGeneralPenitenciaria2026}. Esta población se distribuye en tres subsistemas operativos:

\begin{table}[h!]
\centering
\caption{Distribución de la Población Penal Atendida por Subsistema y Sexo}
\label{tab:poblacion_penal}
\vspace{0.2cm}
\small
\begin{tabular}{lrrr}
\toprule
\textbf{Subsistema Operativo} & \textbf{Hombres} & \textbf{Mujeres} & \textbf{Total Atendidos} \\
\midrule
Subsistema Cerrado (Recintos Penitenciarios) & 61.405 & 5.615 & 67.020 \\
Subsistema Abierto (Penas Sustitutivas) & 62.704 & 10.683 & 73.387 \\
Subsistema Postpenitenciario & 18.766 & 2.870 & 21.636 \\
\midrule
\textbf{Total General} & \textbf{142.875} & \textbf{19.168} & \textbf{162.043} \\
\bottomrule
\end{tabular}
\end{table}

\subsubsection{Radiografía de las Personas Privadas de Libertad (Régimen Cerrado)}
En el régimen cerrado, la población efectivamente recluida supera las 65.000 personas \parencite{CaracterizacionPersonasPrivadas2026}. Este volumen de reclusión ha llevado a las autoridades a alertar sobre niveles críticos de hacinamiento, con un promedio nacional de ocupación que ronda el 130\% y unidades penales de alta densidad que alcanzan hasta un 200\% de su capacidad nominal \parencite{EstadisticaGeneralPenitenciaria2026}.

El perfil demográfico y procesal de las personas privadas de libertad en recintos cerrados presenta los siguientes rasgos estructurales \parencite{CaracterizacionPersonasPrivadas2026}:
\begin{itemize}
    \item \textbf{Calidad Procesal:} El 65,6\% de la población cumple condena efectiva, mientras que el 34,4\% se encuentra en calidad de imputada o procesada.
    \item \textbf{Nacionalidad:} El 83,8\% es de nacionalidad chilena y el 16,2\% corresponde a población extranjera (con predominio de ciudadanos de origen venezolano).
    \item \textbf{Edad y Etnicidad:} El tramo etario con mayor frecuencia se sitúa entre los 25 y los 29 años. En términos de pueblos originarios, la etnia Mapuche representa la mayor proporción de reclusos pertenecientes a personas indígenas.
    \item \textbf{Nivel Educacional y Reincidencia:} El 86\% de las personas en régimen cerrado no cuenta con la enseñanza media completa y la edad promedio bordea los 35 años. La tasa de reincidencia delictiva estimada alcanza un 43,2\%, concentrándose principalmente en delitos contra la propiedad \parencite{CompendiosEstadisticos2026}.
    \item \textbf{Eficiencia Presupuestaria:} El costo fiscal de mantener a una persona privada de libertad en el sistema cerrado asciende aproximadamente a \$800.000 CLP mensuales, mientras que los programas de intermediación y acompañamiento laboral en el medio libre (sistema abierto) requieren una inversión cinco veces menor por usuario \parencite{CompendiosEstadisticos2026}.
\end{itemize}

\subsubsection{Reinserción Social Juvenil y Tendencias de Criminalidad}
En el ámbito infractor juvenil, la implementación del Servicio Nacional de Reinserción Social Juvenil (SNRSJ) reportó durante el primer semestre de 2025 la atención de 4.422 jóvenes a nivel nacional, de los cuales el 93\% corresponde a hombres y el 93\% es de nacionalidad chilena \parencite{Ley21527CreaElServicioNacionalDeReinsercionJuvenil}. Este modelo se fundamenta en programas especializados de intervención psicosocial y mecanismos de compra pública a organismos externos acreditados.

Por su parte, el contexto nacional de criminalidad analizado a través del Centro de Estudios y Análisis del Delito (CEAD) muestra que en 2025 se registraron 1.158 víctimas consumadas de homicidio, consolidando una tasa de 5,7 por cada 100.000 habitantes (una baja del 4,1\% respecto a 2024 y el tercer año consecutivo de descenso) \parencite{InformesEstadisticosCEAD2026}. A pesar de esta tendencia a la baja en delitos violentos graves confirmada durante 2026, la Encuesta Nacional Urbana de Seguridad Ciudadana (ENUSC) señala una victimización del 23,5\% de los hogares y un máximo histórico en la percepción de inseguridad ciudadana del 56,6\% \parencite{ENUSC2025}. Esta brecha entre victimización real y percepción amplifica el estigma social hacia las personas que han cumplido condena.

\subsection{Problematización y Justificación del Estudio}
Los estudios tradicionales de mercado orientados a la inclusión social suelen sufrir de un sesgo crítico: la \textbf{deseabilidad social declarativa}. Cuando se le pregunta directamente a un decisor gerencial si contrataría a una persona con antecedentes o si discrimina a una mujer, la respuesta suele alinearse con la corrección política, ocultando las tasas reales de veto implícito.

Asimismo, existe una desconexión entre la oferta de programas sociales y la estructura de incentivos y riesgos de las empresas según su tamaño (micro, pyme, grande, corporación) y sector económico. Es fundamental precisar empíricamente qué atributos de la oferta de JxR (acompañamiento, nivel técnico, certificaciones CET/OTEC, reducción de riesgo) mitigan eficazmente la resistencia directiva.